% Generated draft: requires mathematical and source review.
\documentclass[11pt]{article}
\usepackage[margin=25mm]{geometry}
\usepackage{fontspec}
\setmainfont{Latin Modern Roman}
\usepackage{amsmath,amssymb,mathtools,unicode-math}
\setmathfont{Latin Modern Math}
\usepackage{enumitem,needspace,xurl,hyperref}
\setlength{\parindent}{0pt}
\setlength{\parskip}{5pt}
\setlength{\emergencystretch}{3em}
\newcommand{\N}{\mathbb{N}}
\newcommand{\Z}{\mathbb{Z}}
\newcommand{\Q}{\mathbb{Q}}
\newcommand{\R}{\mathbb{R}}
\newcommand{\C}{\mathbb{C}}
\newcommand{\F}{\mathbb{F}}
\newcommand{\A}{\mathbb{A}}
\newcommand{\PP}{\mathbb{P}}
\newcommand{\E}{\mathbb{E}}
\newcommand{\eps}{\varepsilon}
\newcommand{\dd}{\,\mathrm{d}}
\begin{document}
\section*{UnsolvedMath ID: 20002441}
\textbf{Source problem code:} catalogue entry
\subsection{Definitions and mathematical statement}
The recovered target is Y. Li's question about stability of finite-depth pure-gravity solitary water waves near their critical speed. The subsequent S. Benzoni-Gavage ``Shocks with Capillarity'' equations and spectral discussion belong to the next presenter, not to this water-wave problem. The preceding companion record contains the water-wave setup, although its displayed nonlinear equations are themselves extraction-damaged.

With depth normalized to one and gravitational acceleration $g>0$, let $(\eta_c,\Phi_c)$ denote a solitary traveling surface elevation and surface velocity potential, depending on $x-ct$ and approaching the rest state at infinity. The flat-surface Dirichlet-to-Neumann multiplier is $q\tanh q$, giving critical long-wave speed
\[
 c_* =\sqrt g.
\]
The question asks how to determine stability of these waves for $c$ near $c_*$. The supplied aid distinguishes longitudinal and transverse linear asymptotic and spectral stability from nonlinear stability of the full free-boundary system. It reports longitudinal linear results of Pego--Sun (2016) for small solitary waves and transverse linear results of Rousset--Sun (2025) for small line solitary waves, while leaving nonlinear full-system stability open. The relevant functional spaces, weights, perturbation classes and exact theorem hypotheses are not included in the record; no unrestricted linear or nonlinear theorem is asserted here.

The aid separately provides an exact flat-state dispersion benchmark, not a stability proof for the nonlinear wave. For longitudinal and transverse frequencies $k,\ell$, set
\[
 q=\sqrt{k^2+\ell^2},\qquad
 \omega(k,\ell)=\sqrt{gq\tanh q}.
\]
Use the right-going sector
\[
 k=\varepsilon K,\quad\ell=\varepsilon^2L,\quad
 c=\frac{\sqrt g}{\sqrt{1-\varepsilon^2}},\quad
 0<\kappa\le K\le M,\quad |L|\le M.
\]
For fixed $\kappa,M$, as $\varepsilon\downarrow0$, the reported expansion is uniform on this sector:
\[
\begin{split}
 \frac{\omega(k,\ell)-ck}{\sqrt g}
 ={}&\varepsilon^3
    \left(\frac{L^2}{2K}-\frac K2-\frac{K^3}{6}\right)\\
 &+\varepsilon^5
    \left(-\frac{L^4}{8K^3}-\frac{KL^2}{4}
          +\frac{19K^5}{360}-\frac{3K}{8}\right)
   +O_{\kappa,M}(\varepsilon^7).
\end{split}
\]
Its leading neutral curve is
\[
 L^2=K^2+\frac{K^4}{3}.
\]
On compact subsets of the stated sector separated from that curve, the moving-frame frequency has the sign of the leading displayed coefficient for all sufficiently small $\varepsilon$. This records the aid's KP-II sign check; it neither classifies the spectrum linearized about a solitary wave nor establishes nonlinear stability. The lower bound $K\ge\kappa>0$ is essential to this uniform expansion and is not silently extended to $K=0$.
\subsection{Short English statement}
Determine near-critical stability of finite-depth gravity solitary waves. The supplied assessment separates published small-wave linear results from the remaining nonlinear question and gives a uniform two-term KP-II dispersion check; the adjacent capillarity-shock problem is separate.
\subsection{Sources}
\begin{itemize}
\item Y. Li, ``Stability of travelling waves of the full water wave problem near the critical case,'' American Institute of Mathematics workshop \emph{Stability criteria for multi-dimensional waves and patterns}, \url{https://aimath.org/WWN/multidimwaves/multidimwaves.pdf}.
\item Local UnsolvedMath registry, permanent ID 20002441, source code \texttt{catalogue entry}, canonical location \texttt{aim-pdes-notes.json}, \texttt{notes[52]}, record dated 14 August 2026. Its verification identifies a crossed presenter boundary; the original literature field is empty.
\item Companion permanent ID 20002440, \texttt{AIM-PDES-0052}, \texttt{notes[51]}, contains Y. Li's heading, the water-wave setup and the flat multiplier $k\tanh k$. Companion permanent ID 20002442, \texttt{AIM-PDES-0054}, \texttt{notes[53]}, identifies the following Benzoni-Gavage material as a separate capillarity problem and dates the workshop to 16--20 May 2005.
\item Supplied research assessment cites Pego--Sun (2016) and Rousset--Sun (2025), without titles, theorem numbers or full bibliographic details. Research attempt 1 is partial result, problem status at run partially solved, rigor labelled proved, novelty confidence low; its dispersion lemma is machine-generated research aid, not peer reviewed, requiring independent verification before citation.
\end{itemize}
\end{document}
